% Scritto da pitch/figure.py a partire da grafica/dati/poteri.json: non modificarlo a mano.
\newcommand{\potereUno}{Passaparola}
\newcommand{\effettoUno}{Pesca 1 carta Ambito.}
\newcommand{\breveUno}{pesca 1 Ambito.}
\newcommand{\chiUno}{08 · 10 · 12 · 13 · 15 · 18 · 23 · 24 · 33 · 34 · 35 · 49 · 55 · 57 · 58}
\newcommand{\potereDue}{Colloquio}
\newcommand{\effettoDue}{Prendi una Persona dal mercato e tienila da parte, scoperta: è prenotata per te. Il mercato si ripristina.}
\newcommand{\breveDue}{prenota una Persona del mercato.}
\newcommand{\chiDue}{04 · 07 · 09 · 14 · 19 · 28 · 38 · 42 · 44 · 46 · 47 · 52 · 53 · 54 · 60}
\newcommand{\potereTre}{Corso intensivo}
\newcommand{\effettoTre}{Gioca subito fino a 2 carte Ambito dalla mano come Formazioni: per aprire pile nuove o per migliorarne di tue.}
\newcommand{\breveTre}{gioca fino a 2 Ambiti come formazioni.}
\newcommand{\chiTre}{05 · 16 · 17 · 25 · 26 · 32 · 37 · 39 · 43 · 45 · 48 · 50 · 51 · 56 · 59}
\newcommand{\potereQuattro}{Chiamata diretta}
\newcommand{\effettoQuattro}{Metti subito un'altra Persona, dal mercato o tra quelle prenotate, su una tua pila libera, ignorando una delle sue tre ambizioni. Il suo potere non si attiva.}
\newcommand{\breveQuattro}{piazza un'altra Persona ignorando un'ambizione (senza il suo potere).}
\newcommand{\chiQuattro}{01 · 02 · 03 · 06 · 11 · 20 · 21 · 22 · 27 · 29 · 30 · 31 · 36 · 40 · 41}
